% ---
% title: Claims on two policies
% tags: finance, joint-discrete, marginal, joint-distribution-function
% ---

\begin{problem}
	\begin{context}
		Many customers hold several policies with the same insurer, for
		example a car policy and a home policy. The insurer pays for every
		claim, so it matters how the claims on the different policies occur
		together: a customer who claims on one policy may also be more likely
		to claim on the other. The joint distribution of the
		claim counts describes this.
	\end{context}
	A customer has a car policy and a home policy with the same insurer. Let
	$X$ be the number of claims on the car policy in a year ($0$ or $1$), and
	$Y$ the number of claims on the home policy in the same year ($0$, $1$ or
	$2$). The joint probability mass function $p(x,y)=P(X=x,Y=y)$ is
	\begin{center}
		\begin{tabular}{c|ccc}
			          & $y=0$  & $y=1$  & $y=2$  \\ \hline
			$x=0$     & $0.50$ & $0.15$ & $0.05$ \\
			$x=1$     & $0.15$ & $0.10$ & $0.05$
		\end{tabular}
	\end{center}
	(a) Find the marginal probability mass functions $p_X$ and $p_Y$.
	(b) Find the probability that the customer makes at least two claims in
	total, $P(X+Y\ge2)$.
	(c) Find the joint distribution function $F(a,b)=P(X\le a,Y\le b)$ at the
	points $(0,0)$, $(0,1)$, $(1,0)$ and $(1,1)$.
\end{problem}

\begin{solution}
	\idea Every probability about $X$ and $Y$ is a sum of cells of the table:
	find the cells where the event happens and add them. The marginal pmfs are
	the row sums and the column sums. $F(a,b)$ is the sum of the cells with
	$x\le a$ and $y\le b$: a ``lower left'' block of the table.

	\step{Check the table.} The six cells are all the possible pairs
	$(x,y)$, so they must add up to $1$:
	\[ 0.50+0.15+0.05+0.15+0.10+0.05=1. \]

	\step{(a)} The event $\{X=0\}$ happens together with exactly one of
	$\{Y=0\}$, $\{Y=1\}$, $\{Y=2\}$. These three are mutually exclusive, so
	\[ p_X(0)=P(X=0)=p(0,0)+p(0,1)+p(0,2)=0.50+0.15+0.05=0.70. \]
	In words: $p_X(0)$ is the sum of the row $x=0$. In the same way, the row
	$x=1$ gives
	\[ p_X(1)=p(1,0)+p(1,1)+p(1,2)=0.15+0.10+0.05=0.30. \]
	For $Y$ we add the columns instead:
	\[
		\begin{aligned}
			p_Y(0) & =p(0,0)+p(1,0)=0.50+0.15=0.65, \\
			p_Y(1) & =p(0,1)+p(1,1)=0.15+0.10=0.25, \\
			p_Y(2) & =p(0,2)+p(1,2)=0.05+0.05=0.10.
		\end{aligned}
	\]
	It is convenient to write these sums into the margins of the table (this
	is where the name ``marginal'' comes from):
	\begin{center}
		\begin{tabular}{c|ccc|c}
			          & $y=0$  & $y=1$  & $y=2$  & $p_X(x)$ \\ \hline
			$x=0$     & $0.50$ & $0.15$ & $0.05$ & $0.70$   \\
			$x=1$     & $0.15$ & $0.10$ & $0.05$ & $0.30$   \\ \hline
			$p_Y(y)$  & $0.65$ & $0.25$ & $0.10$ & $1$
		\end{tabular}
	\end{center}
	Check: $0.70+0.30=1$ and $0.65+0.25+0.10=1$. So the car policy has no
	claim in 70\% of the years, and the home policy has no claim in 65\% of
	the years.

	\step{(b)} List the pairs $(x,y)$ with $x+y\ge2$. With $x=0$ we need
	$y=2$; with $x=1$ we need $y\ge1$, so $y=1$ or $y=2$. The three cells
	$(0,2)$, $(1,1)$, $(1,2)$ are different outcomes, so they are mutually
	exclusive, and their probabilities add up:
	\[ P(X+Y\ge2)=p(0,2)+p(1,1)+p(1,2)=0.05+0.10+0.05=0.20. \]
	Check with the complement: $X+Y\le1$ means $(0,0)$, $(0,1)$ or $(1,0)$,
	with probability $0.50+0.15+0.15=0.80=1-0.20$.

	\step{(c)} $F(a,b)$ adds the cells with $x\le a$ \emph{and} $y\le b$.
	\begin{itemize}
		\item $F(0,0)=P(X\le0,Y\le0)$. Both counts are at least $0$, so this is
		      just the one cell $(0,0)$: $F(0,0)=p(0,0)=0.50$.
		\item $F(0,1)$: $x=0$ and $y\in\{0,1\}$, two cells:
		      $F(0,1)=p(0,0)+p(0,1)=0.50+0.15=0.65$.
		\item $F(1,0)$: $x\in\{0,1\}$ and $y=0$, two cells:
		      $F(1,0)=p(0,0)+p(1,0)=0.50+0.15=0.65$.
		\item $F(1,1)$: $x\in\{0,1\}$ and $y\in\{0,1\}$, four cells:
		      \[ F(1,1)=p(0,0)+p(0,1)+p(1,0)+p(1,1)=0.50+0.15+0.15+0.10=0.90. \]
	\end{itemize}
	A faster way for $F(1,1)$: the only cells left out are those with $y=2$,
	so $F(1,1)=1-P(Y=2)=1-0.10=0.90$.

	\begin{answer}
		(a) ${p_X(0)=0.70}$, ${p_X(1)=0.30}$; ${p_Y(0)=0.65}$, ${p_Y(1)=0.25}$,
		${p_Y(2)=0.10}$ \quad
		(b)~${P(X+Y\ge2)=0.20}$ \quad
		(c) ${F(0,0)=0.50}$, ${F(0,1)=0.65}$, ${F(1,0)=0.65}$, ${F(1,1)=0.90}$
	\end{answer}
	\begin{remark}
		$F(0,1)$ and $F(1,0)$ are equal here, but they are sums of different
		cells: $F(0,1)$ uses $(0,0)$ and $(0,1)$, $F(1,0)$ uses $(0,0)$ and
		$(1,0)$. They agree only because $p(0,1)=p(1,0)=0.15$.
		Since $X\le1$ always, $F(1,b)=P(Y\le b)$ is the distribution function
		of $Y$ alone: for example $F(1,0)=p_Y(0)=0.65$.
		The claims do tend to occur together:
		$P(X=1,Y=1)=0.10$ is larger than $p_X(1)\,p_Y(1)=0.30\cdot0.25=0.075$,
		so a car claim makes a home claim more likely here, and $X$ and $Y$ are
		not independent.
		\textbf{Common mistake:} computing $F(1,1)$ as $P(X=1,Y=1)=0.10$. The
		joint distribution function uses ``$\le$'', so it collects a whole
		block of cells, not a single cell.
	\end{remark}
\end{solution}
